\clearpage\pagestyle{otMain}\setlength{\otBlocEspacement}{2.5mm}\fontsize{11}{15}\selectfont\corrigesChapitre{3}{Multiples, diviseurs et partages}
\begin{blocLivre}{cor-3-1}\begin{corrige}{1}{Multiples à repérer}\demarche{Compter de 7 en 7 ou utiliser un produit.}
$7, 14, 21, 28, 35$. Et $7\times12=84$ : $\res{\text{oui}}$, 84 est un multiple de 7.
\end{corrige}\end{blocLivre}
\begin{blocLivre}{cor-3-2}\begin{corrige}{2}{Les critères des unités}\demarche{Lire uniquement le chiffre des unités.}
Multiples de 2 : $\res{230\text{ et }1\,308}$. Multiples de 5 : $\res{230,\ 345,\ 95}$. Multiples de 10 : $\res{230}$ seulement.
\end{corrige}\end{blocLivre}
\begin{blocLivre}{cor-3-3}\begin{corrige}{3}{Diviseur ou pas ?}\demarche{Chercher un produit exact ou poser la division.}
$8\times12=96$ : $\res{\text{oui}}$. $93=7\times13+2$ : $\res{\text{non}}$, le reste vaut 2.
\end{corrige}\end{blocLivre}
\begin{blocLivre}{cor-3-4}\begin{corrige}{4}{Tous les diviseurs}\demarche{Balayer les facteurs dans l’ordre croissant.}
$1\times24$, $2\times12$, $3\times8$, $4\times6$ : la liste est $\res{1, 2, 3, 4, 6, 8, 12, 24}$.
\end{corrige}\end{blocLivre}
\begin{blocLivre}{cor-3-5}\begin{corrige}{5}{Pairs et impairs}\demarche{Raisonner sur les paquets de deux.}
Pair plus pair : $\res{\text{paire}}$. Impair plus impair : $\res{\text{paire}}$. Pair plus impair : $\res{\text{impaire}}$. Les paquets de deux expliquent ces règles.
\end{corrige}\end{blocLivre}
\begin{blocLivre}{cor-3-6}\begin{corrige}{6}{Diviseurs de 30 et 18}\demarche{Comparer les deux listes complètes.}
Diviseurs de 30 : $1,2,3,5,6,10,15,30$. Diviseurs de 18 : $1,2,3,6,9,18$. Les diviseurs communs sont $\res{1, 2, 3, 6}$.
\end{corrige}\end{blocLivre}
\begin{blocLivre}{cor-3-7}\begin{corrige}{7}{Multiples communs}\demarche{Comparer les deux listes de multiples.}
Multiples de 6 : 6, 12, 18, 24, 30, 36, 42, 48. Multiples de 8 : 8, 16, 24, 32, 40, 48. Les communs sont 24 et 48 ; le plus petit est $\res{24}$.
\end{corrige}\end{blocLivre}
\begin{blocLivre}{cor-3-8}\begin{corrige}{8}{Un chiffre à choisir}\demarche{Appliquer les critères au chiffre des unités.}
Multiple de 5 : le nombre doit finir par 0 ou 5, donc $\square=\res{0\text{ ou }5}$. Multiple de 10 : $\square=\res0$ seulement.
\end{corrige}\end{blocLivre}
\begin{blocLivre}{cor-3-9}\begin{corrige}{9}{Le reste raconte}\demarche{Relier reste nul et diviseur.}
Le reste vaut 5, pas 0 : 6 $\res{\text{ne divise pas}}$ 125. Le multiple suivant est $6\times21=\res{126}$.
\end{corrige}\end{blocLivre}
\begin{blocLivre}{cor-3-10}\begin{corrige}{10}{Vrai ou faux ?}\demarche{Chercher un contre-exemple pour réfuter.}
La première est $\res{\text{vraie}}$ : un nombre fini par 0 est aussi fini par 0 ou 5 ; $10=5\times2$, donc dix paquets de 5. La seconde est $\res{\text{fausse}}$ : 15 est multiple de 5 mais pas de 10.
\end{corrige}\end{blocLivre}
\begin{blocLivre}{cor-3-11}\begin{corrige}{11}{Les bouquets}\demarche{Chercher les diviseurs communs des deux effectifs.}
Le nombre de bouquets doit diviser 24 et 18 : $\res{1, 2, 3\text{ ou }6}$ bouquets. Avec 6 bouquets : $24:6=\res4$ roses et $18:6=\res3$ tulipes chacun.
\end{corrige}\end{blocLivre}
\begin{blocLivre}{cor-3-12}\begin{corrige}{12}{Les deux phares}\demarche{Chercher le plus petit multiple commun.}
Multiples de 12 : 12, 24, 36… Multiples de 9 : 9, 18, 27, 36… Le premier commun est $\res{36}$ secondes.
\end{corrige}\end{blocLivre}
\begin{blocLivre}{cor-3-13}\begin{corrige}{13}{Les rangées du verger}\demarche{Tester chaque diviseur par une division.}
$90=5\times18$ : $\res{18}$ rangées exactes. $90=6\times15$ : $\res{15}$ rangées exactes. $90=7\times12+6$ : impossible sans reste, il resterait $\res6$ arbres.
\end{corrige}\end{blocLivre}
\begin{blocLivre}{cor-3-14}\begin{corrige}{14}{Le nombre secret}\demarche{Lister puis filtrer les candidats.}
Diviseurs de 24 : 1, 2, 3, 4, 6, 8, 12, 24. Les pairs multiples de 3 sont 6, 12 et 24 : trois candidats. La dernière contrainte, « plus petit que 10 », ne laisse que $\res6$. Chaque indice d’une devinette élimine des candidats, jusqu’à n’en garder qu’un.
\end{corrige}\end{blocLivre}
\begin{blocLivre}{cor-3-15}\begin{corrige}{15}{Toujours vrai ?}\demarche{Dépasser l’exemple par l’argument des paquets.}
Exemple : $10+25+40=\res{75}=5\times15$. Chaque terme est fait de paquets de cinq ; leur réunion aussi : l’affirmation est $\res{\text{toujours vraie}}$. L’argument vaut pour tous les choix, pas seulement l’exemple.
\end{corrige}\end{blocLivre}
